\section{DISCUSIÓN}

El valor de la vigilancia genómica en aguas residuales está en convertir la señal molecular en información sobre la circulación viral de una población. Los resultados de este trabajo confirman que el esquema aplicado en la PTAR Quitumbe recupera material viral de los tres agentes patógenos seleccionados con un rendimiento de caracterización condicionado por la carga viral y por la conservación de las muestras. Esta sección interpreta cada conjunto de resultados frente a la literatura revisada y al estudio antecedente local de SARS-CoV-2 \cite{MejiaCalle_2024}.

\subsection{Detección molecular en una matriz ambiental compleja}

La detección de los tres virus en distintas etapas de la investigación confirma que el muestreo pasivo, así como el flujo de concentración y extracción, recuperan material genético viral del influente de la planta. Este resultado refleja el principio establecido desde los primeros reportes de SARS-CoV-2 en aguas residuales, donde la detección de ARN viral acompañó la circulación comunitaria aun con vigilancia clínica limitada \cite{Ahmed_2020,Medema_2020}. La capacidad de detectar SARS-CoV-2, un enterovirus no polio y RSV con un mismo procedimiento respalda la viabilidad operacional de un esquema multipatógeno en un único sitio centinela, una ventaja descrita para la vigilancia basada en aguas residuales en contextos con recursos limitados \cite{Parkins_2023,Bubba_2023}.

\subsection{Carga viral, valores de Cq y rendimiento de la secuenciación}

Las variaciones observadas en los valores de Cq en las muestras explican el rendimiento desigual de la caracterización genómica. Los análisis con linaje que se logró determinar, concentraron los valores de Cq más bajos del gen N1, mientras que las muestras con Cq tardíos generaron consenso con un alto porcentaje de nucleótidos indeterminados e incluso, no se secuenciaron. Esta descripción coincide con lo reportado previamente en la PTAR Quitumbe, donde las muestras con Cq tardíos fueron menos eficientes para la secuenciación \cite{MejiaCalle_2024}, y con la observación general de que las cargas virales bajas y el ARN fragmentado limitan la recuperación genómica en aguas residuales \cite{Parkins_2023,Child_2023}. El mecanismo de esa pérdida reside en la propia PCR por amplicones, donde la formación de dímeros entre cebadores es la causa principal del sesgo de cobertura y se acentúa a medida que baja la carga viral \cite{Itokawa_2020}. Un Cq tardío no invalida la detección; confirma la presencia del virus, permite seguir sus tendencias de circulación y fija el umbral práctico a partir del cual la secuenciación deja de ser informativa; en aguas residuales, solo las muestras con Cq inferior a 33 rindieron genomas de consenso completos \cite{Crits_Christoph_2021}. En Dhaka, un programa de vigilancia por Nanopore en un país de ingresos bajos seleccionó por su valor de Cq, las muestras enviadas a secuenciación y recuperó genomas con cobertura del 94,2 al 99,8 \% a partir de muestras con Cq de ORF1ab entre 28,8 y 32,1 \cite{Haque_2023}. La amplificación de un solo gen diana (N1 o N2) refuerza esta interpretación, dado que la señal cercana al límite de detección se vuelve intermitente, ciertas veces se detecta un gen, otras veces el otro, y no necesariamente ambos.

\subsection{Linajes de SARS-CoV-2 y contexto epidemiológico}

Los linajes de SARS-CoV-2 identificados entre las semanas 3 y 9 de 2025 corresponden por completo a descendientes de JN.1: JN.1, LZ.5 y LU.2.1.1 se ubican en el clado 24A, definido por JN.1, y LF.7.1.3 en el clado 24H, también derivado de JN.1, según la asignación de Nextclade. Este predominio concuerda con lo observado en la región andina, donde el linaje 23I (JN.1) fue dominante a lo largo de 2024 \cite{Bruno_2025}. La detección de JN.1 en aguas residuales de Quito reproduce el hallazgo del antecedente local, que reportó este linaje antes de su identificación clínica \cite{MejiaCalle_2024}, y confirma la capacidad de la secuenciación de aguas residuales para identificar variantes regionalmente prevalentes cuando la cobertura es suficiente \cite{Crits_Christoph_2021}.

\citeA{Bruno_2025} reunieron en GISAID 12\,619 genomas ecuatorianos de SARS-CoV-2 del periodo 2020--2024, el 1,15\,\% de los casos confirmados del país. El aporte anual cayó de 5\,102 genomas en 2022 a 744 en 2024. Esa serie no llega a 2025, el año del muestreo, y deja los linajes del influente sin un referente clínico ecuatoriano contemporáneo \cite{Bruno_2025}.

\subsection{\textit{Enterovirus}}

En dos muestras prospectivas (semanas 45 y 46), la PCR anidada recuperó enterovirus, identificado por BLAST como Coxsackievirus A, con amplicones de la región VP4 en lugar de la VP1 prevista. La vigilancia ambiental de poliovirus emplea la recuperación de enterovirus no poliovirus como control de proceso, y los protocolos de la OMS esperan aislarlos en al menos el 30 \% de las muestras puntuales concentradas y el 10 \% de las obtenidas con instrumentos de muestreo \cite{OReilly_2020,Bubba_2023}. Con el dispositivo pasivo, este estudio los recuperó en 2 de 22 muestras (9 \%), justo por debajo del umbral del 10 \%, lo que confirma que el muestreo y la extracción captaron ARN enteroviral pero refleja la menor sensibilidad del muestreo pasivo frente al puntual concentrado \cite{Rozanec_2026}. El sesgo hacia VP4 apunta a una especificidad de los cebadores de VP1 menor que la esperada en esta matriz, que el protocolo deberá corregir; aun así, el esquema detectó enterovirus que la vigilancia de parálisis flácida aguda no registra \cite{Asghar_2014}.

Ninguna de las 22 muestras rindió poliovirus. El esquema ecuatoriano mantiene la vacuna oral bOPV —a los 6 y 18 meses y a los 5 años, además de dos dosis de polio inactivada en la vacuna hexavalente— \cite{MSP_2023}, y la excreción de poliovirus Sabin en el alcantarillado sería esperable tras la vacunación \cite{MasLago_2003}. Que las 22 muestras no lo recuperaran apunta a la baja sensibilidad del muestreo pasivo, que aquí rindió enterovirus en apenas el 9\,\%, antes que a una circulación nula. El poliovirus sigue siendo un blanco pertinente. En Suiza, con esquema exclusivo de vacuna inactivada, el alcantarillado de Zúrich rindió poliovirus tipo Sabin \cite{Zurbriggen_2008}, y en Eslovaquia décadas de vigilancia en aguas residuales aislaron poliovirus vacunal pero nunca salvaje \cite{Kissova_2022,Asghar_2014}. El flujo aplicado aquí se ajusta a ese objetivo. \citeA{Seo_2025} recuperaron en Canadá, en 2022, un poliovirus derivado de vacuna de tipo 2 cuya región VP1 tuvo un 99,4 \% de identidad con aislados ambientales de Londres y un 99,0 \% con el caso clínico de Nueva York, sin poliomielitis asociada y compatible con la excreción de un viajero; el cultivo de las muestras conservadas falló y la secuencia se obtuvo por PCR anidada y secuenciación Nanopore, el mismo flujo de este estudio \cite{Seo_2025}. La secuenciación de VP1 distingue las cepas vacunales de las salvajes \cite{Klapsa_2022}, y clasificaría una eventual detección en la PTAR Quitumbe.

\subsection{RSV}

La señal de RSV en cuatro muestras sin asignación de subgrupo reproduce el comportamiento esperado del virus en aguas residuales. La recuperación de genomas completos exige cargas altas —con cobertura buena por encima de 10\,000 copias/mL— que las muestras ambientales de baja carga difícilmente alcanzan \cite{Le_2025}. A ese límite se suma que RSV es un virus envuelto, cuya recuperación desde las aguas residuales es intrínsecamente baja; cuando el ARN va empaquetado en partículas envueltas, la eficiencia de recuperación cae del 71 \% al 17 \% frente al ARN desnudo \cite{ToribioAvedillo_2024}. El subgrupo agrava el problema, ya que RSV-A domina la positividad ambiental sobre RSV-B —32,6 \% frente a 2,4 \% en ocho plantas de tratamiento del sureste de Alemania—, y su recuperación depende sobre todo de los métodos de concentración y extracción \cite{Geissler_2025}. \citeA{Hughes_2022} correlacionaron el ARN de RSV en sólidos sedimentados con la positividad clínica por RT-qPCR, con coeficientes de Kendall de 0,65 a 0,77, sin secuenciar el virus. Las cuatro muestras positivas de este estudio reprodujeron esa detección cuantitativa, pero el análisis con EPI2ME no alineó ninguna con la referencia de RSV y el subgrupo quedó sin asignar.

\subsection{Conservación de muestras retrospectivas como factor crítico}

Las ocho muestras retrospectivas, conservadas en DNA/RNA Shield 2X y en PBS 1X, no rindieron secuencia de poliovirus ni de RSV; todas las secuenciaciones provinieron de muestras prospectivas frescas. \citeA{Barbe_2022} midieron ese efecto al secuenciar SARS-CoV-2 por Nanopore, donde la cobertura del genoma tuvo una mediana del 94\,\% partiendo de ARN fresco y del 88\,\% con el extracto de ARN congelado, sin diferencia significativa entre ambos, pero cayó al 55\,\% a partir de agua residual congelada. Nuestras muestras retrospectivas se conservaron sobre la matriz completa, la vía que allí más redujo la cobertura, y su rendimiento fue nulo. \citeA{MejiaCalle_2024} había ligado el rendimiento del muestreo pasivo en la PTAR Quitumbe a la integridad del ARN y a la eficiencia de recuperación, y \citeA{Parkins_2023} documentan que la congelación y el calentamiento de las aguas residuales degradan sus ácidos nucleicos virales.

El almacenamiento perturba incluso la cuantificación. \citeA{Zambrana_2024} remidieron RSV y el marcador fecal PMMoV en sólidos de aguas residuales conservados a $-80$\,°C hasta 700 días y con dos ciclos de congelación-descongelación, y obtuvieron valores significativamente distintos de los de las muestras frescas (Kruskal-Wallis, RSV $P=1{,}6\times10^{-3}$; PMMoV $P=1{,}7\times10^{-10}$). El RSV se movió menos de un orden de magnitud, con una razón almacenado/fresco de 1,29, mientras que el PMMoV descendió a 0,40 de su valor fresco. El almacenamiento desplaza así las medidas cuantitativas de forma moderada, pero compromete la secuenciación hasta impedirla, como ocurrió con estas muestras retrospectivas.

